%!TeX root=csp-proof.tex
\section{The hard half}\label{sec:hardness}

If no weak near-unanimity operation preserves $\Gamma$, then $\CSP(\Gamma)$ is
\textup{NP}-complete. This half is older, shorter, and better understood than
the tractable half, and it factors cleanly into algebra and complexity.

\subsection{The chain}

\begin{description}\tightlist
\item[\textup{(1)}] \emph{Reduction to a core, and to an idempotent algebra.}
  One may assume $\Gamma$ is a rigid core containing all singleton unary
  relations, so that $\Pol(\Gamma)$ is idempotent, at the cost of a
  logspace reduction \cite[``Cores and Idempotent Reducts'']{BradyNotes}.
  This is Caveat~\ref{haz:core}.
\item[\textup{(2)}] \emph{Maróti--McKenzie.} A finite idempotent algebra has a
  Taylor term if and only if it has a weak near-unanimity term
  \cite{MarotiMcKenzie}.
\item[\textup{(3)}] \emph{No Taylor term $\Rightarrow$ two disjoint
  subalgebras.} If a finite idempotent $\mathbf A$ has no Taylor term then
  there are nonempty $B,C\le\mathbf A$ with $B\cap C=\varnothing$ and
  \[
    N=(B\cup C)^{3}\setminus(B^{3}\cup C^{3})\;\le\;\mathbf A^{3}
  \]
  \cite[Cor.~1.5.11]{BradyNotes}. This replaces the usual formulation --- ``the
  variety generated contains a two-element algebra whose only operations are
  projections, and from there $\Gamma$ pp-interprets every finite language'' ---
  whose ``and'' hides an \textup{HSP} representation in a finite power and its
  translation into an interpretation. The cited corollary does that work and
  lands on a concrete relation, from which Lemma~\ref{lem:nae-interp} builds
  the interpretation in four lines.
\item[\textup{(3$'$)}] \emph{Inv--Pol.} $N$ is a subuniverse of $\mathbf A^{3}$
  and so is preserved by every polymorphism of $\Gamma$; Theorem~\ref{thm:galois}
  turns that into pp-definability from $\Gamma$. This is the one substantial
  import of the step, and only the hard direction of it is used.
\item[\textup{(4)}] \emph{Bulatov--Jeavons--Krokhin.} A primitive positive
  interpretation of $\Delta$ in $\Gamma$ yields a polynomial-time many-one
  reduction $\CSP(\Delta)\le_{p}\CSP(\Gamma)$ \cite{BJK}.
\item[\textup{(5)}] \emph{Cook--Levin.} $\textsc{Nae-$3$-Sat}$ is
  \textup{NP}-hard.
\end{description}

Steps (1)--(3) are algebra; step (4) is a syntactic construction plus a
complexity-theoretic wrapper; step (5) is pure complexity theory.

\subsection{The layered targets}

\begin{lemma}[The interpretation, explicitly]\label{lem:nae-interp}
Let $\mathbf A$ be an algebra, $B,C\le\mathbf A$ nonempty and disjoint, and
suppose $N=(B\cup C)^{3}\setminus(B^{3}\cup C^{3})$ is a subuniverse of
$\mathbf A^{3}$. Then $\{N\}$ pp-interprets $\{\textsc{Nae}\}$ on $\{0,1\}$,
with $d=1$.
\end{lemma}

\begin{proof}
A triple lies in $N$ exactly when its entries lie in $B\cup C$ and are not all
on the same side of the partition; in particular, for a repeated entry,
\begin{equation}\label{eq:nae-opposite}
  (x,z,z)\in N\iff x,z\in B\cup C\ \text{and $x,z$ lie on opposite sides.}
\end{equation}
Take $F=\proj_{1}(N)$, which is pp-definable from $N$ and equals $B\cup C$:
each $b\in B$ is the first entry of $(b,c,c)\in N$ for any $c\in C$, and
symmetrically. Let $\pi\colon F\to\{0,1\}$ send $B$ to $0$ and $C$ to $1$; it is
surjective because $B$ and $C$ are nonempty. The $\pi$-preimage of
$\textsc{Nae}$ is $N$ itself, by the description above. The $\pi$-preimage of
equality on $\{0,1\}$ is $B^{2}\cup C^{2}$, and by \eqref{eq:nae-opposite}
\[
  B^{2}\cup C^{2}=\{(x,y):\exists z\;N(x,z,z)\wedge N(y,z,z)\} :
\]
a $z$ opposite to both $x$ and $y$ exists exactly when $x$ and $y$ are on the
same side, since there are only two sides and both are nonempty.
\end{proof}

\begin{theorem}[Target \textup{(H0)}]\label{thm:h0}
Let $\mathbf A$ be a finite idempotent algebra with no weak near-unanimity term
operation. Then $\Inv(\mathbf A)$ primitively positively interprets the
constraint language $\{\textsc{Nae}\}$ on $\{0,1\}$, where
$\textsc{Nae}=\{0,1\}^{3}\setminus\{(0,0,0),(1,1,1)\}$.
\end{theorem}

\begin{proof}
By Lemma~\ref{lem:special-wnu}'s source \cite{MarotiMcKenzie}, a finite
idempotent algebra has a Taylor term if and only if it has a weak
near-unanimity term; so $\mathbf A$ has no Taylor term. Step \textup{(3)}
supplies nonempty disjoint $B,C\le\mathbf A$ with
$N=(B\cup C)^{3}\setminus(B^{3}\cup C^{3})\le\mathbf A^{3}$, so
$N\in\Inv(\mathbf A)$, and Lemma~\ref{lem:nae-interp} builds the
interpretation. Every relation it uses is pp-definable from $N$ and hence lies
in $\Inv(\mathbf A)$.
\end{proof}

\begin{theorem}[Target \textup{(H1)}; stated, not proved]\label{thm:h1}
If $\Gamma$ primitively positively interprets $\Delta$ then
$\CSP(\Delta)\le_{p}\CSP(\Gamma)$. Consequently, under \textup{(H0)},
$\CSP(\Gamma)$ is \textup{NP}-hard.
\end{theorem}

\begin{hazard}[Hardness is not completeness]\label{haz:np-membership}
Theorem~\ref{thm:h1} concludes \textup{NP}-\emph{hardness}. The word
\emph{complete} in Theorem~\ref{thm:dichotomy-core} additionally needs
membership, which is Theorem~\ref{thm:np-membership}. Small, but it needs the
cost model, so it belongs with \textup{(H1)} and not with \textup{(H0)}.
\end{hazard}

\begin{theorem}[\textup{NP} membership; stated, not proved]\label{thm:np-membership}
For each finite $\Gamma$, $\CSP(\Gamma)\in\textup{NP}$.
\end{theorem}

\begin{proof}[Proof sketch]
A solution is a certificate: an assignment of a value of $A$ to each variable,
of size $O(v\log|A|)$ for an instance with $v$ variables, hence linear in the
instance. Verifying it evaluates each constraint once, and for fixed $\Gamma$
that is a table lookup of constant cost. Writing this out needs the encoding
and the machine, which is why it is stated with \textup{(H1)} rather than
proved here.
\end{proof}

With Theorems~\ref{thm:h1} and \ref{thm:np-membership} together, the hard half
of Theorem~\ref{thm:dichotomy-core} reads \textup{NP}-\emph{complete} as
claimed.

Theorem~\ref{thm:h0} mentions no machine. Theorem~\ref{thm:h1} is where the
complexity layer enters, and it is small --- the reduction is a syntactic
substitution of formulas --- but it cannot be stated at all without a precise
notion of polynomial-time many-one reduction.

\subsection{Definitions}

\begin{definition}[pp-formula, pp-definability]\label{def:pp}
A \emph{primitive positive formula} over $\Gamma$ is
$\exists y_{1}\dots\exists y_{s}\;\bigwedge_{j}R_{j}(\bar z_{j})$ with
$R_{j}\in\Gamma\cup\{=\}$. A relation is \emph{pp-definable} from $\Gamma$ if
some pp-formula defines it.
\end{definition}

\begin{definition}[pp-interpretation]\label{def:pp-interp}
$\Gamma$ on $A$ \emph{pp-interprets} $\Delta$ on $B$ if there are $d\ge 1$, a
pp-definable $F\subseteq A^{d}$, and a surjection $\pi\colon F\to B$ such that
the $\pi$-preimage of every relation of $\Delta$, and of the equality relation
on $B$, is pp-definable from $\Gamma$.
\end{definition}

\begin{theorem}[Geiger; Bodnarchuk--Kalu\v znin--Kotov--Romov]\label{thm:galois}
For a finite $A$, a relation is pp-definable from $\Gamma$ if and only if it is
preserved by every polymorphism of $\Gamma$.
\end{theorem}

\begin{hazard}[The Galois connection is the expensive import]\label{haz:galois}
Theorem~\ref{thm:galois} is what converts the algebraic hypothesis (``no WNU
polymorphism'') into a syntactic conclusion (``pp-defines a hard relation''),
and it is the only genuinely substantial import of \S\ref{sec:hardness}. The
easy direction is a one-line calculation. The hard direction --- every relation
preserved by all polymorphisms is pp-definable --- goes through the free algebra
on $|A|^{k}$ generators and an indicator-instance construction; it is roughly
four printed pages and is the single largest item in the \textup{(H0)} budget.
\end{hazard}

\begin{remark}[Why the hard half is not the hard half]\label{rem:hardness-easier}
Despite its name, \textup{(H0)} is far smaller than \textup{(T0)}: perhaps
twenty printed pages against a hundred. The asymmetry is intrinsic. Showing a
problem is hard requires exhibiting one construction; showing it is easy
requires an algorithm that works on every instance, and the correctness proof
must survive every configuration the instance can present.
\end{remark}
